\documentclass[10pt,a4paper]{amsart}
\usepackage[margin=19mm]{geometry}
\usepackage{amsmath,amssymb,mathtools}
\usepackage[scr=rsfs,cal=euler]{mathalfa}
\usepackage{xfrac}
\usepackage[unicode,hidelinks,bookmarks=false]{hyperref}
\newcommand{\ie}{i.e.\ }
\newcommand{\cf}{cf.\ }
\newcommand{\resp}{respectively\ }
\newcommand{\Subsection}{\subsection}
\providecommand{\nobreakdash}{\nobreak\hbox{-}}
\setlength{\emergencystretch}{2em}
\multlinegap0pt
\usepackage{tikz}
\usepackage{float}
\usepackage{amssymb}
\usepackage{dsfont}
\usepackage{cancel}
\usepackage{braket}
\newtheorem{lem}{Lemma}
\newcommand		{\CK}			{\mathfrak{c}}	% hat{K} = \CK / |k|
\newcommand		{\Gr}			{G^{\textnormal{r}}}
\newcommand		{\GrE}			{G^{\textnormal{r},E}}
\newcommand		{\Nrm}[2]		{\nrm{#1}_{#2}}
\newcommand		{\R}		{\RR}
\newcommand		{\RR}		{\mathbb R}			% real numbers
\newcommand		{\Rd}		{\R^3}
\newcommand		{\bNrm}[2]		{\bnrm{#1}_{#2}}
\newcommand		{\bangle}[1]	{\lt\langle #1\rt\rangle}
\newcommand		{\bnrm}[1]		{\big\lVert #1\big\rVert}
\newcommand	{\cM}		{\mathcal M}		% Mesures bornées
\renewcommand		{\d}		{\mathop{}\!\mathrm{d}}		% différentielle
\newcommand		{\eps}			{\varepsilon}
\newcommand	{\h}			{h}
\newcommand		{\hatG}			{\widehat{G}}
\newcommand		{\hatGr}		{\hatG^{\textnormal{r}}}
\newcommand		{\hatGrE}		{\hatG^{\textnormal{r},E}}
\newcommand		{\intd}			{\int_{\Rd}}
\newcommand		{\lt}			{\left}				%
\newcommand		{\n}[1]			{\lt\lvert{#1}\rt\rvert}	% |.|
\newcommand		{\nrm}[1]		{\lt\lVert{#1}\rt\rVert}		% ||.||
\newcommand		{\rt}			{\right}			%
\newcommand		{\sNrm}[2]		{\snrm{#1}_{#2}}
\newcommand		{\snrm}[1]		{\lVert #1\rVert}
\newcommand		{\weight}[1]	{\bangle{#1}}	% <x>
\title{\Large Quantum Landau Damping with Coulomb Repulsion in the Whole Space}
\author{Dongfen Bian, Emmanuel Grenier, Laurent Lafl\`eche, Quoc Hung Nguyen}
\date{Source: arXiv:2609.02040v1 (CC BY 4.0)}
\begin{document}
\maketitle

\noindent\emph{Source and licence.} \Large Quantum Landau Damping with Coulomb Repulsion in the Whole Space. Version arXiv:2609.02040v1 (CC BY 4.0), distributed by arXiv under the Creative Commons Attribution 4.0 International licence, http://creativecommons.org/licenses/by/4.0/, as recorded in the arXiv licence metadata for this exact version and shown on the abstract page https://arxiv.org/abs/2609.02040v1. Full licence notice: \texttt{licenses/arxiv-2609.02040-CC-BY-4.0.txt}.\par\smallskip

\noindent\textit{Editorial scope.} Counted unit: the article's lem lem:decay\_Gr (source0.tex lines 1797--1819). The article's complete original proof of it is delivered with it (source0.tex lines 1821--1930, one \texttt{proof} environment of 110 lines). No mathematics is written, repaired or re-derived: the statement and the proof are the article's own lines, in the article's own notation, and no retained line is substituted or re-flowed.  Retained uncounted as the source-local context the proof needs: lines 1335--1344.  Nothing was omitted from any retained span, so the package carries no omission record.  No retained line contains a citation, so the wrapper carries no bibliography and no bibliography entry.  No retained line contains a figure environment, an image-inclusion command or drawing code, so no image is delivered and the package declares no asset dependency.  Editorial formatting only, in addition: an AMS \texttt{amsart} preamble that restates the article's own declarations and macros used by the retained lines, namely the environments \texttt{lem} on separate counters, together with the standard packages those lines need.  The retained environments print in this document as Lemma 1 in order of appearance, whereas the article prints the same items with the numbering of its own full document.  The complete native source of the article as distributed in the arXiv e-print for this exact version is bundled unchanged as \texttt{sources/209165/source0.tex}.
		\begin{align}\label{eq:Gr-symbol-derivatives}
			\n{\partial_k^\alpha\hatGr_\h(t,k)}
			&\leq C \, \frac{\n{k}^{3-\n{\alpha}}}{\weight{k}^4}
			\weight{\n{k}t}^{\n{\alpha}}e^{-\sigma\n{k}t},
			\\\label{eq:Gr-difference-symbol-derivatives}
			\n{\partial_k^\alpha
			\bigl(\hatGr_\h-\hatGr_0\bigr)(t,k)}
			&\leq C \, h^2\, \frac{\n{k}^{5-\n{\alpha}}}{\weight{k}^4}
			\weight{\n{k}t}^{\n{\alpha}} e^{-\sigma\n{k}t} .
		\end{align}

	\begin{lem}\label{lem:decay_Gr}
		For any $h\in[0,1]$ and $t>0$, it holds that
        %{\color{blue}(The short times blowups in $L^1$ can probably be deleted by considering a sharper interpolation estimate with $|x|^3$ instead of %\eqref{eq:weighted-L2-to-L1}, if all the previous estimates work for $3+\eps$ derivatives)}
		\begin{align*}
			\Nrm{\Gr_\h(t,\cdot)}{L^\infty} &\lesssim \frac{1}{t^2\weight{t}^4} \,  ,
			&
			\Nrm{\Gr_\h(t,\cdot)}{L^1} &\lesssim \frac{1}{t^{1/8}\weight{t}^{3-1/8}} \,  ,
			\\
			\Nrm{\GrE_\h(t,\cdot)}{L^\infty} &\lesssim \frac{1}{t\weight{t}^4} \,  ,
			&
			\Nrm{\GrE_\h(t,\cdot)}{L^1} &\lesssim \frac{1}{\weight{t}^{2}} \, . 
		\end{align*}
		Moreover, we have
		\begin{align*}
			\Nrm{\Gr_\h(t,\cdot) - \Gr_0(t,\cdot)}{L^\infty} &\lesssim \frac{h^2}{t^4\weight{t}^4} \, ,
			&
			\Nrm{\Gr_\h(t,\cdot) - \Gr_0(t,\cdot)}{L^1} &\lesssim \frac{h^{2}}{t\weight{t}^{4}} % Or \frac{h^{2}}{t^{1/2}\weight{t}^{5-1/2}} \, ,
			\\
			\Nrm{\GrE_\h(t,\cdot) - \GrE_0(t,\cdot)}{L^\infty} &\lesssim \frac{h^2}{t^3\weight{t}^4} \, ,
			&
			\Nrm{\GrE_\h(t,\cdot) - \GrE_0(t,\cdot)}{L^1} &\lesssim \frac{h^{2}}{t^{3/8}\weight{t}^{4-3/8}} \, .
		\end{align*}
	\end{lem}

	\begin{proof}
		We first record two elementary scaling estimates. For $m>-3$ and $t>0$, we have
		\begin{equation}\label{eq:radial-L1-scaling}
			\intd \n{k}^m e^{-\sigma \n{k}t} \d k = \tfrac{4\pi \Gamma(m+3)}{\sigma^{m+3}}\, t^{-m-3} \, .
		\end{equation}
		Similarly, if $m-a \geq - 3/2$, then
		\begin{equation}\label{eq:radial-L2-scaling}
			\bNrm{\n{k}^{m-a}\weight{\n{k}t}^{a}e^{-\sigma \n{k} t}}{L^2} = C\, t^{-m+a-3/2}
		\end{equation}
		for some finite constant $C>0$ depending on $a$, $k$ and $\sigma$. On the other hand, when $t\to 0$, then by dominated convergence, for $p\geq 1$
		\begin{align*} %\label{eq:radial-L2-scaling_t_small}
			&\Nrm{\weight{k}^{-4}\n{k}^{m-a}\weight{\n{k}t}^{a}e^{-\sigma \n{k} t}}{L^p} \to C &&\text{ if } -\frac{3}{p}<m-a<4-\frac{3}{p} \, ,
            \\
            &\Nrm{\weight{k}^{-4}\n{k}^{m-a}\weight{\n{k}t}^{a}e^{-\sigma \n{k} t}}{L^p} \sim C\,t^{4-3/p+a-m} &&\text{ if } m-a>4-\frac{3}{p} \, .
		\end{align*}

		Let
		\begin{equation*}
			\mathcal M_E(k) := -\widehat{\nabla K}(k) = -2i\pi \CK \, \frac{k}{\n{k}^2} \, ,
			\qquad
			\hatGrE_\h(t,k)=\mathcal M_E(k)\hatGr_\h(t,k) \, .
		\end{equation*}
		Since $\mathcal M_E$ is homogeneous of degree $-1$, for any multi-index $\beta$ one has
		\begin{equation*}
			\n{\partial_k^\beta\cM_E(k)} \leq C_\beta \n{k}^{-1-\n{\beta}}  .
		\end{equation*}
		Leibniz' rule and the differentiated bounds~\eqref{eq:Gr-symbol-derivatives} and~\eqref{eq:Gr-difference-symbol-derivatives} therefore imply, for $\n{\alpha}\leq 2$,
		\begin{align}\label{eq:GrE-symbol-derivatives}
			\n{\partial_k^\alpha \hatGrE_\h(t,k)}
			&\lesssim \frac{\n{k}^{2-\n{\alpha}}}{\weight{k}^4}
			\weight{\n{k}t}^{\n{\alpha}}e^{-\sigma \n{k}t},
			\\\label{eq:GrE-difference-symbol-derivatives}
			\n{\partial_k^\alpha
			\bigl(\hatGrE_\h - \hatGrE_0\bigr)(t,k)}
			&\lesssim h^2 \, \frac{\n{k}^{4-\n{\alpha}}}{\weight{k}^4}
			\weight{\n{k}t}^{\n{\alpha}}e^{-\sigma \n{k}t}.
		\end{align}
		Indeed, every term in the first Leibniz sum is bounded by
		\begin{equation*}
			\n{k}^{-1-|\beta|}
			\frac{\n{k}^{3-|\alpha-\beta|}}{\weight{k}^4}
			\weight{\n{k}t}^{|\alpha-\beta|}e^{-\sigma \n{k}t}
			\leq \frac{\n{k}^{2-\n{\alpha}}}{\weight{k}^4}
			\weight{\n{k}t}^{\n{\alpha}}e^{-\sigma \n{k}t},
		\end{equation*}
		and the estimate on the difference is identical with $\n{k}^3$ replaced by $h^2\n{k}^5$. Together with inequalities~\eqref{eq:Gr-symbol-derivatives} and~\eqref{eq:Gr-difference-symbol-derivatives}, these estimates show that the derivatives up to order two of all four symbols are locally square integrable at $k=0$. 

		We now prove the $L^\infty$ bounds. The Fourier inversion formula gives 
		\begin{align*}
			\Nrm{\Gr_\h(t,\cdot)}{L^\infty} \leq\sNrm{\hatGr_\h(t,\cdot)}{L^1} &\lesssim t^{-6},
			\\
			\Nrm{\GrE_\h(t,\cdot)}{L^\infty} \leq\sNrm{\widehat{\GrE_\h}(t,\cdot)}{L^1} &\lesssim t^{-5},
			\\
			\Nrm{\Gr_\h(t,\cdot)-\Gr_0(t,\cdot)}{L^\infty} &\lesssim \h^2\, t^{-8},
			\\
			\Nrm{\GrE_\h(t,\cdot)-\GrE_0(t,\cdot)}{L^\infty} &\lesssim \h^2 \, t^{-7}.
		\end{align*}
		It remains to establish the $L^1$ bounds. We shall use the following weighted interpolation inequality in dimension three
		\begin{equation}\label{eq:weighted-L2-to-L1}
			\Nrm{g}{L^1} \lesssim \Nrm{g}{L^2}^{1/4} \Nrm{\n{x}^2g}{L^2}^{3/4} ,
		\end{equation}
		which follows for instance by writing for $R>0$, by the Cauchy--Schwarz inequality,
		\begin{align*}
			\Nrm{g}{L^1} = \int_{B_R} \n{g} \, dx + \int_{\Rd\setminus B_R} \n{g} \, dx 
			\lesssim R^{3/2}\Nrm{g}{L^2} + R^{-1/2}\Nrm{\n{x}^2g}{L^2}
		\end{align*}
		and optimizing in $R$. By Plancherel's theorem
		\begin{equation}\label{eq:weighted-plancherel}
			\Nrm{\n{x}^2g}{L^2}
			=\frac{1}{4\pi^2} \Nrm{\Delta_k\widehat g}{L^2}
			\lesssim
			\sum_{\n{\alpha}=2}
			\Nrm{\partial_k^\alpha\widehat g}{L^2}.
		\end{equation}
		Applying Formula~\eqref{eq:radial-L2-scaling} to \eqref{eq:Gr-symbol-derivatives} and~\eqref{eq:GrE-symbol-derivatives} yields
		\begin{align*}
			\Nrm{\Gr_\h(t)}{L^2} &\lesssim t^{-9/2},
			&\Nrm{\n{x}^2\Gr_\h(t)}{L^2}
			&\lesssim t^{-5/2},\\
			\Nrm{\GrE_\h(t)}{L^2}
			&\lesssim t^{-7/2},
			&\Nrm{\n{x}^2\GrE_\h(t)}{L^2}
			&\lesssim t^{-3/2}.
		\end{align*}
		Hence the interpolation inequality~\eqref{eq:weighted-L2-to-L1} gives
		\begin{equation*}
			\Nrm{\Gr_\h(t)}{L^1}\lesssim t^{-3},
			\qquad
			\Nrm{\GrE_\h(t)}{L^1}\lesssim t^{-2}.
		\end{equation*}
		For $0<t\leq1$, retain the factor $\weight{k}^{-4}$ in~\eqref{eq:GrE-symbol-derivatives} instead of discarding it in the scaling estimate. The cases $\n{\alpha}=0$ and $\n{\alpha}=2$ then give
		\begin{equation*}
			\Nrm{\GrE_\h(t)}{L^2} + \Nrm{\n{x}^2\GrE_\h(t)}{L^2}\leq C \, .
		\end{equation*}
		Indeed, $\weight{t k}^2 = 1+t^2\n{k}^2$, and both resulting radial integrals are bounded uniformly for $0<t\leq1$ after the change of variables $r=t\n{k}$ in the term containing $t^2$. Thus \eqref{eq:weighted-L2-to-L1} also gives $\sNrm{\GrE_\h(t)}{L^1}\leq C$ on this interval. Similarly, applying Formula~\eqref{eq:radial-L2-scaling} and Inequality~\eqref{eq:weighted-plancherel} to the two differences of symbols gives
		\begin{align*}
			\Nrm{\Gr_\h(t)-\Gr_0(t)}{L^2} &\lesssim \h^2 \, t^{-13/2},
			&\Nrm{\n{x}^2(\Gr_\h-\Gr_0)(t)}{L^2} &\lesssim \h^2 \, t^{-9/2},
			\\
			\Nrm{\GrE_\h(t)-\GrE_0(t)}{L^2} &\lesssim \h^2 \, t^{-11/2},
			&\Nrm{\n{x}^2(\GrE_\h-\GrE_0)(t)}{L^2} &\lesssim \h^2 \, t^{-7/2}.
		\end{align*}
		A last use of \eqref{eq:weighted-L2-to-L1} gives
		\begin{equation*}
			\Nrm{\Gr_\h(t)-\Gr_0(t)}{L^1} \lesssim \h^2 \, t^{-5},
			\qquad
			\Nrm{\GrE_\h(t)-\GrE_0(t)}{L^1} \lesssim \h^2 \, t^{-4}.
		\end{equation*}
		This completes the proof.
	\end{proof}

\end{document}
